\section{Related Work and Motivation}
This section reviews prior research and motivates our study.

\subsection{Software Evolution}

\emph{The one constant in software is change}~\cite{sokolowski2022change}. Software evolution is commonly categorized as adaptive, corrective, perfective, and preventive~\cite{zhang2024developers}. A substantial body of research has examined software maintenance and refactoring activities~\cite{baqais2020automatic,bavota2015experimental,martini2015investigating,vassallo2019large}. For example, Bavota \textit{et al.}~\cite{bavota2015experimental} found that 42\% of refactoring activities are driven by code smells. Another major research direction focuses on understanding change- and bug-prone software components~\cite{arvanitou2017method,chowdhury2022empirical,chowdhury:2022,Shihab:2012,PASCARELLA:2020,chowdhury2024method,Mo:2022}. Such studies aim to identify defects early and reduce maintenance cost, effort, and reputational risk.

In contrast to production code, relatively few studies focused exclusively on test code evolution. A common assumption is that test code evolves primarily in response to changes in production code~\cite{sun_revisiting_2023,hu2023identify}, which may render tests obsolete or redundant. Consequently, much of the research has centered on production--test co-evolution~\cite{zaidman2011studying,sun_revisiting_2023,wang2021understanding}, test obsolescence~\cite{hao2013bug}, test deletion~\cite{bhatta2025understanding}, and automated test repair~\cite{yaraghi2025automated}.

Zaidman \textit{et al.}~\cite{zaidman2011studying} investigated whether production and test code evolve synchronously or in separate phases, observing evidence of both patterns. Wang \textit{et al.}~\cite{wang2021understanding} identified common factors (e.g., modified language constructs) that indicate whether a test should change following a production code modification. Similarly, Hu \textit{et al.}~\cite{hu2023identify} proposed a transformer-based approach for repairing test cases affected by production code changes. However, the authors relied on name matching to map production and test methods, an approach that suffers from low recall (as we show in \textbf{RQ2}), causing many test cases requiring repair to be missed. Another limitation in production--test co-evolution research concerns how co-evolution is identified. Most studies assume that when a test artifact and its associated production artifact are modified together, or within a short temporal window, the pair represents an instance of co-evolution. However, this assumption was found to produce a large number of false positives~\cite{sun_revisiting_2023}.

Pinto \textit{et al.}~\cite{pinto2012understanding} identified test repair as a primary driver of test evolution and showed that assertion-focused repair techniques are often insufficient for handling complex real-world repairs. However, without any specialized tool, tracking test cases across commits proved challenging because many test methods in their dataset were renamed or moved over time. Bhatta \textit{et al.}~\cite{bhatta2025understanding} examined why and how test methods are deleted in real-world Java applications. They found that most deleted tests are obsolete, while some others are redundant. Interestingly, deleting approximately 20\% of redundant tests reduced test suite effectiveness by lowering coverage or mutation scores. Again, due to the lack of a specialized tool, the authors relied on heuristics, \emph{RefactoringMiner}~\cite{tsantalis2020refactoringminer}, and extensive manual verification for building their test deletion dataset. This multi-dimensional process not only limited scalability but may also have affected accuracy, since \emph{RefactoringMiner} exhibits relatively low recall (54.2\%) for detecting extract \& move refactorings.

In a different study, Spadini \textit{et al.}~\cite{Spadini:2018} found that test methods containing test smells are more change- and bug-prone than those without smells. To reconstruct test method histories, the authors assumed that two methods were identical if they shared more than 95\% cosine similarity when an exact signature match could not be found. However, the accuracy of this assumption was not validated against ground truth, unlike more recent method-level history tracking techniques~\cite{grund:2021,jodavi_accurate_2022,islam_historyfinder_2025}. Furthermore, the study did not investigate what proportion of test method changes were induced by changes in their corresponding production methods, a question that requires accurate method-level production-to-test mapping.

\emph{Taken together, these studies suggest that accurately tracking test method histories and establishing production-to-test mappings at the method level are crucial for an accurate understanding of test code evolution.}

\subsection{Method-level history generation}

Early tools such as Historage~\cite{hata_2011} and FinerGit~\cite{HIGO:2020} relied on Git-based mechanisms and required expensive offline preprocessing. CodeShovel~\cite{grund:2021} addressed this limitation by generating method histories on demand using a string-similarity-based tracking algorithm. However, a follow-up study by Jodavi \textit{et al.}~\cite{jodavi_accurate_2022} identified limitations in both CodeShovel's algorithm and the oracle used for its evaluation. They subsequently proposed CodeTracker~\cite{jodavi_accurate_2022}, which leverages \emph{RefactoringMiner}~\cite{tsantalis2020refactoringminer} to track method histories and demonstrated improved accuracy on a corrected benchmark. More recently, Islam \textit{et al.}~\cite{islam_historyfinder_2025} revisited both method-history generation and its evaluation. They showed that existing tools were assessed using inaccurate ground-truth oracles and constructed a new, systematically developed benchmark for evaluation. Using this benchmark, they reevaluated prior approaches and introduced \emph{HistoryFinder}, which achieved the best overall trade-off between tracking accuracy and execution time.

\emph{Unfortunately, these high-performing tools—CodeShovel, CodeTracker, and HistoryFinder—were developed and evaluated exclusively on production methods. As test code differs substantially from production code in many aspects~\cite{van2021comparative,horikawa2025does,zaidman2011studying}, we do not know whether the string-based matching heuristics and thresholds~\cite{grund:2021, islam_historyfinder_2025} that perform well for production methods generalize to test methods without recalibration.}

% These tools were primarily designed and evaluated for production source code. However, their effectiveness on test code cannot be assumed directly. Test methods often differ from production methods in structure and intent: they commonly contain setup logic, assertions, mocks, fixtures, framework annotations, parameterized executions, and behavior-oriented names. Test code may also contain shorter methods, repeated assertion patterns, duplicated scenarios, and helper methods whose evolution is tied to testing intent rather than production API design. As a result, similarity thresholds, refactoring heuristics, and code metrics calibrated for production methods may not transfer cleanly to test methods. Therefore, method-level history generation for test code requires separate evaluation, and potentially recalibration, before these tools can be reliably used in test-code evolution studies.

\subsection{Method-level production-to-test mapping}
Linking production and test code is essential for regression test selection, traceability, co-evolution analysis, and test maintenance~\cite{white_establishing_2020,white_tctracer_2022,sun_method-level_2024,liu_cote_2025}. However, establishing these links is challenging, particularly at the method level~\cite{white_establishing_2020}. Test methods often invoke helper, setup, assertion, and utility methods in addition to the production methods they validate, causing static and dynamic analyses to incorrectly identify incidental calls as tested methods~\cite{white_tctracer_2022,sun_method-level_2024}.

Van Rompaey and Demeyer~\cite{van_rompaey_establishing_2009} found that simple heuristics such as \emph{Naming Conventions (NC)}, which matches test and production method names, and \emph{Last Call Before Assert (LCBA)}, which identifies the last production method invoked before an assertion, achieve high accuracy. However, these approaches are less effective when naming conventions are inconsistent or when a test exercises multiple production methods~\cite{liu_cote_2025,white_tctracer_2022}. For example, LCBA can incorrectly identify helper methods as the tested method due to assertion proximity. White \textit{et al.}~\cite{white_establishing_2020} evaluated several additional linking techniques, including \emph{Naming Conventions Contains (NCC)}, which relaxes exact name matching; \emph{Longest Common Subsequence} using both names for normalization (LCS-B) or only the production name (LCS-U); and \emph{Levenshtein distance (Leven)}. They also explored \emph{Tarantula}~\cite{jones2002visualization}, which applies fault-localization principles to identify tested methods, and \emph{Term Frequency–Inverse Document Frequency (TFIDF)}, which treats tests as documents and production methods as terms. Using a manually constructed oracle from three open-source Java projects, they found that a combined approach incorporating all techniques achieved the highest accuracy. In a later study using an expanded benchmark, the authors showed that incorporating dynamic traces substantially improves the accuracy of naming-based approaches~\cite{white_tctracer_2022}.

Using dynamic traces, however, is challenging because projects must compile and tests must be executable, limiting scalability for mining software repositories (MSR) studies that analyze large numbers of projects~\cite{sun_method-level_2024}. To address this limitation, Sun \textit{et al.}~\cite{sun_method-level_2024} proposed \emph{TestLinker}, a purely static approach that leverages semantic correlation learning and a pretrained code model. \emph{TestLinker} was evaluated on an extended benchmark that added two projects to the dataset used by White \textit{et al.}~\cite{white_tctracer_2022}.

\emph{Prior research showed that the relative effectiveness of competing techniques often depends on the benchmark used for evaluation~\cite{lee2018bench4bl,chowdhury2024method}, and we observe a similar phenomenon here. For example, the NC heuristic achieved precisions of 100\%, 100\%, and 93\% in~\cite{van_rompaey_establishing_2009,white_tctracer_2022,sun_method-level_2024}, respectively, while the corresponding recalls were 100\%, 11\%, and 12\%. Such variation makes it difficult to determine the most effective approach. Therefore, before applying production-to-test mapping to the 110 projects studied in this paper, we systematically reevaluate existing techniques using a new benchmark of 20 open-source projects, substantially larger than the benchmarks of 3 to 6 projects used in prior studies~\cite{white_establishing_2020,white_tctracer_2022,sun_method-level_2024}.}
% Recent work has moved beyond heuristics by learning semantic
% relationships between tests and focal methods. Sun \textit{et al.}~\cite{sun_method-level_2024}
% proposed TestLinker, which frames method-level traceability as semantic
% correlation learning and uses a pretrained code model to rank candidate focal
% methods. This reduces dependence on hand-crafted lexical heuristics, but the
% approach still relies on mapping rules to connect
% inferred method names to concrete implementation methods.

% Method-level test-to-production linking identifies the production method or
% methods that a test method exercises or validates. T
% Linking production and test code is essential for regression test selection, traceability, co-evolution analysis, and test maintenance~\cite{white_establishing_2020,white_tctracer_2022,sun_method-level_2024,liu_cote_2025}. However, establishing these links is challenging, more so at the method level than class level~\cite{white_establishing_2020}. Test methods often invoke helper, setup, assertion, and utility methods in addition to the production methods they are intended to validate. As a result, static or dynamic call extraction can mistakenly identify incidental calls as tested production methods~\cite{white_tctracer_2022,sun_method-level_2024}. 

% Van Rompaey and Demeyer~\cite{van_rompaey_establishing_2009} found that simple \emph{Naming Conventions (NC)} where the test method is named by pre or post-pending the name of the production method to the word \emph{test}, and \emph{Call Before Assert (LCBA)} where the most recently called production method before assertion is considered as the actual production method that is tested, are both highly accurate. However, these approaches may not perform well when a project does not follow naming conventions very consistently or when one test exercises multiple production methods~\cite{liu_cote_2025, white_tctracer_2022}. For example, LCBA that relies on assertion proximity can misidentify helper calls as the identified production method. In addition to \emph{Naming Conventions (NC)}, White \textit{et al.}~\cite{white_establishing_2020} explored other linking strategies including \emph{Naming Conventions Contains (NCC)}, which relaxes NC by replacing exact
% name equality with a substring match; \emph{Longest Common Subsequence} between test and production names using both
% names (LCS-B) for normalization and using only the production name for normalization (LCS-U); and Levenshtein distance (Leven), which scores edit-distance similarity. The authors also leveraged approaches like \emph{Tarantula}~\cite{jones2002visualization} to use fault localization techniques for mapping to the correct production method, and \emph{Term Frequency–Inverse Document Frequency (TFIDF)} such that tests are treated as documents and production methods are treated as the terms. By using an oracle of links constructed from three open-source Java projects, the authors found that the \emph{combined} techniques that include all the aforementioned approaches show the best accuracy.  

% TCTracer selected NC and LCBA as established baselines and formulated additional
% method-level linking strategies
% \cite{white_establishing_2020,white_tctracer_2022}. Its string-based strategies
% include Naming Conventions Contains (NCC), which relaxes NC by replacing exact
% name equality with a substring match; Longest Common Subsequence using both
% names (LCS-B) and using the unit under test (LCS-U), which score ordered name
% overlap; and Levenshtein distance (Leven), which scores edit-distance
% similarity. These strategies broaden exact-name matching, but they still depend
% on meaningful name overlap. TCTracer also includes statistical call-based
% strategies, such as Tarantula and TFIDF, which emphasize distinctive
% test--method call relationships without directly relying on names. However,
% such strategies may still rank helper or utility calls highly when those calls
% are specific to a test. TCTracer's Combined score mitigates some individual
% weaknesses by aggregating multiple signals, although it can still inherit noise
% from its components. Together, these studies show that links inferred from
% names, raw call relationships, or co-change signals can miss relevant tested
% methods or include incidental methods that are not the behavior under test.

% Recent work has moved beyond heuristics by learning semantic
% relationships between tests and focal methods. Sun \textit{et al.}~\cite{sun_method-level_2024}
% proposed TestLinker, which frames method-level traceability as semantic
% correlation learning and uses a pretrained code model to rank candidate focal
% methods. This reduces dependence on hand-crafted lexical heuristics, but the
% approach still relies on mapping rules to connect
% inferred method names to concrete implementation methods.

% Despite this progress, the resolution from candidate links to concrete
% implementation methods remains a potential source of error. Dynamic approaches,
% including TCTracer, use execution traces or coverage to observe which methods run during a test
% \cite{gergely_differences_2019,white_tctracer_2022}, whereas TestLinker maps the
% inferred focal method name to a concrete production declaration using heuristic
% mapping rules \cite{sun_method-level_2024}. In contrast, our pipeline is static:
% it first uses Java Symbol Solver to resolve method calls to their concrete
% implementation methods and to filter irrelevant calls, and uses heuristic
% mapping only as a fallback when symbol resolution is insufficient. Evaluation
% also depends on the quality of the benchmark ground truth. Prior ground truths
% may miss or mislabel method-level links, so an apparent false positive may
% sometimes be a correct link that is absent from, or incorrectly labeled. RQ2 therefore compares the linking strategies under a common setting and
% reports inaccuracies observed in prior ground-truth links.
%\subsection{Test smells}

% In this section, we describe prior studies on Software evolution, method-level change tracking, test-to-production mapping, and code smells. 
% \subsection{Software evolution} \emph{The one constant in software is change}~\cite{sokolowski2022change}. Software evolution is generally classified as adaptive, corrective, perfective, and preventive~\cite{zhang2024developers}. Refactoring for code maintenance has been widely studied~\cite {baqais2020automatic, bavota2015experimental,martini2015investigating,vassallo2019large}. For example, Bavota \textit{et al.}~\cite{bavota2015experimental} studied the correlation between software quality and refactoring and observed that 42\% of refactoring activities are induced by code smells. Another common area of research is understanding change and bug-prone code components~\cite {arvanitou2017method,chowdhury2022empirical,chowdhury:2022, Shihab:2012,PASCARELLA:2020, chowdhury2024method, Mo:2022}. In bug-proneness research, for example, the idea is to predict bug early to reduce maintenance cost, effort, and reputation. 

% Unfortunately, compared to production code evolution, only a few studies have focused solely on test code evolution. There is a prevailing assumption in the community that test code evolution happens mainly in response to changes in production code~\cite{sun_revisiting_2023,hu2023identify}, which makes the test code obsolete or redundant. Based on this assumption, research has focused on product-test co-evolution~\cite{zaidman2011studying,sun_revisiting_2023,wang2021understanding}, test obsolescence~\cite{hao2013bug}, test deletion~\cite{bhatta2025understanding}, and automatic test repair~\cite{yaraghi2025automated}. Zaidman \textit{et al.}~\cite{zaidman2011studying} investigated whether co-evolution happens synchronously or in different phases. The observation was mixed: sometimes phased and sometimes synchronous evolution was observed. 
% Wang \textit{et al.}~\cite{wang2021understanding} observed that there are common factors (e.g., modified language constructs) that are good indicators of whether a test code should change based on a change in production code. Based on this observation, the authors developed a machine learning-based approach to predict obsolete tests that need repairs. A similar work was done by Hu \textit{et al.}~\cite{hu2023identify} who proposed a transformer-based approach to identify and repair test cases in response to a corresponding production code change. One particular strength of this work is that the recommendation is made at method level instead of class level. However, for mapping between production and test methods, the authors relied on name matching, which has very low recall as we show in our RQ2. Another common pitfall in studies of production–test co-evolution lies in how co-evolution is identified. Most approaches assume that when a test artifact and its associated production artifact are modified together, or within a short temporal window, the pair constitutes an instance of production–test co-evolution. This approach, however, suffers from high false positives~\cite{sun_revisiting_2023}. 

% Pinto \textit{et al.}~\cite{pinto2012understanding} investigated why test suites evolve and found that test repair is one of the prevailing reasons for test code evolution, justifying the importance of developing effective test repair techniques. The authors showed that simple assertion-focused test repair techniques are less applicable, as many test repairs in real-world settings are much more complex. However, the authors needed to track test cases across commits, which was challenging because many test methods were renamed or moved in their dataset. In another similar study, Bhatta \textit{et al.}~\cite{bhatta2025understanding} studied why and how test methods are deleted in real-world Java applications. The authors found that a major fraction of deleted tests are
% obsolete tests while some are redundant. Surprisingly, the deletion of 20\% of redundant tests reduces test suite effectiveness as they reduce coverage or mutation scores. To capture the samples of deleted tests, however, the authors relied on heuristics, RefactoringMiner~\cite{tsantalis2020refactoringminer}, and laborious manual verification. The difficulty of applying this multi-dimensional approach not only impacted the scalability of the study, but also could impact the accuracy, as RefactoringMiner shows the lowest recall (54.2) among all other transformations, in detecting Extract & Move Method. 

% In a different study, Spadini \textit{et al.}~\cite{Spadini:2018} observed that test methods with test smells are more change- and bug-prone than test methods without smells. To collect the change history of the studied test methods, when a method with the same signature is not found, the authors relied on the assumption that if two methods have more than 95\% cosine similarity, they are the same methods. Unfortunately, the accuracy of this assumption in capturing the whole change history of a method is not validated as was done for recent method-level change history detectors~\cite{grund:2021, jodavi_accurate_2022, islam_historyfinder_2025}. Also, the authors did not investigate what portions of those test method changes were induced by their corresponding source method changes, which require method-level production-to-test mapping. 

% \emph{Considering the previous studies, we argue that to properly understand test code evolution the importance of accurately tracking test code history and production-to-test mapping at the method level is paramount.}
